\chapter{\conclusionsname}

El desarrollo de las comunicaciones cuánticas y de la ciberseguridad en la era post-cuántica pone de manifiesto que la computación cuántica no solo representa una promesa tecnológica, sino un desafío para las infraestructuras digitales actuales. Desde enfoques distintos pero complementarios, ambas áreas convergen en un mismo objetivo, garantizar la seguridad, eficiencia y resiliencia de los sistemas de comunicación en un mundo cada vez más interconectado.

Los hallazgos permiten extraer varias reflexiones generales. En primer lugar, la computación cuántica se presenta como motor de innovación y, al mismo tiempo, como fuente de vulnerabilidad, obligando a replantear tanto las infraestructuras de comunicación como los modelos de seguridad que las protegen. En segundo lugar, la estandarización y la interoperabilidad resultan indispensables para una adopción global, sin ellas, los avances tan solo quedarían como proyectos experimentales. En tercer lugar, la colaboración entre expertos de diferentes disciplinas será un factor decisivo para acelerar la transición hacia un ecosistema cuántico verdaderamente confiable.

De cara al futuro, todo apunta a que la próxima década estará marcada por avances significativos en las tecnologías cuánticas, con potencial para consolidarse de manera similar a lo ocurrido con la inteligencia artificial, que pasó en pocos años de ser una promesa lejana a estar presente hoy en día en sectores tan diversos como la salud, la industria o la gestión empresarial, entre otros. Las tecnologías cuánticas también transformarán la manera en que se entiende la información, el conocimiento y la interacción con el mundo. Ofrecen nuevas formas de procesar datos, interpretar fenómenos complejos y abordar problemas que desafían las estructuras matemáticas tradicionales, produciendo una revolución en la realidad tecnológica y científica. Imaginar el impacto que podría suponer la combinación de ambas disciplinas, la inteligencia artificial y las tecnologías cuánticas.

Si bien en las comunicaciones cuánticas ya se han conseguido desplieges operativos experimentales, aún quedan muchos desafíos por superar antes de pensar en la construcción de un internet cuántico global. Esto muestra la importancia de trabajos como este, los cuales contribuyen signifivativamente al avance de la investigación organizando el conocimiento existente y facilitando el desarrollo hacia nuevos paradigmas.

Por todo ello, estas tecnologías podrían llegar a ser una de las transformaciones más disruptivas de la historia, con un potencial que va mucho más allá de la informática.
